\begin{table}[H]
\centering
\small
\setlength{\tabcolsep}{3.5pt}
\caption{\textbf{Comparable external model-family audit on public cohorts.}
Each entry is an equal-weight mean across nine runs (three cohorts $\times$
three seeds) on the corresponding cohort-specific frozen 4,096-gene panel and
condition-disjoint
holdout construction. Per-seed test-cell counts match across the audited
model outputs. Cosine is higher-is-better; MSE and MAE are lower-is-better.
The GEARS rows are the verified CUDA runs in Table~\ref{tab:external-gears};
the remaining rows are deterministic or fitted baselines evaluated with the
same split construction. These summaries are conventional response-prediction
controls and are not pooled with the Frangieh reliability-transport claim.}
\label{tab:external-model-summary}
\begin{tabular}{@{}lrrrrr@{}}
\toprule
Model family & Runs & Cohorts & Cosine & MSE & MAE \\
\midrule
Ridge regression & 9 & 3 & 0.233 & 0.01887 & 0.09970 \\
Perturbation mean $\Delta$ & 9 & 3 & 0.202 & 0.01705 & 0.09548 \\
Global mean $\Delta$ & 9 & 3 & 0.185 & 0.01737 & 0.09624 \\
Context $k$NN $\Delta$ (fallback) & 9 & 3 & 0.185 & 0.01737 & 0.09624 \\
GEARS & 9 & 3 & 0.122 & 0.02073 & 0.10533 \\
Control mean & 9 & 3 & 0.000 & 0.01823 & 0.09852 \\
\bottomrule
\end{tabular}
\par\noindent
\raggedright\footnotesize Equal-run means are not test-cell-weighted; the K562
GWPS cohort contributes only 64 non-control test cells per three-run cohort.
Under the perturbation-disjoint holdout, the $k$NN control falls back to the
global mean for unseen test perturbations; it is retained as a protocol
control, not as evidence of context-neighbour transfer. Full run-level values
and output paths are in
\texttt{results/tables/external\_condition\_holdout\_shared\_panel\_metrics.csv}.
\end{table}
